\section{Introducción}
La gestión de talento humano (GTH) impacta productividad, clima y cumplimiento. En Las Brisas, procesos manuales y documentos dispersos aumentan tiempos de respuesta, errores y falta de trazabilidad. La digitalización de procesos críticos —asistencias, permisos y vacaciones, contratos, formación y evaluaciones— es prioritaria. Este artículo sintetiza la propuesta técnica, los métodos de validación y los resultados esperados para una plataforma de GTH adaptada al contexto de Las Brisas. La propuesta trasciende una ERS tradicional al incluir benchmarking, análisis costo–beneficio, métricas de adopción, riesgos y gobernanza de datos, con el fin de sustentar la inversión estratégica con beneficios medibles y sostenibles.

\section{Marco teórico}
La propuesta se fundamenta en cuatro bases principales: Análisis de requisitos, que permite definir claramente las necesidades del sistema y asegurar que se puedan verificar más adelante. Modelos de calidad de software, que plantean que una buena solución debe ser eficiente, segura, compatible y fácil de usar. Diseño centrado en las personas, con el que se espera mejorar la experiencia y satisfacción de los usuarios mediante pruebas y mejoras progresivas. Protección y manejo de datos, que incluye reglas de seguridad, control de accesos, registros de auditoría y transparencia en el uso de la información. Estos elementos orientan el desarrollo de una solución modular, segura y fácil de usar, adaptada a las necesidades de gestión de personal en Las Brisas.